\section{Momentum}

\SourcePage{62}{65}
Consider a closed system of particles in no external field. Since all
positions of the system as a whole are equivalent, its Hamiltonian is unchanged
by an arbitrary parallel displacement. It is enough to require this for every
infinitesimal displacement; it then holds for finite displacements as well.

Under an infinitesimal displacement $\delta\mathbf r$, every radius vector
$\mathbf r_a$ acquires the same increment. Thus
$\mathbf r_a\to\mathbf r_a+\delta\mathbf r$, and
\[
 \psi(\mathbf r_1+\delta\mathbf r,\mathbf r_2+\delta\mathbf r,\ldots)
 =\left(1+\delta\mathbf r\cdot\sum_a\boldsymbol\nabla_a\right)
 \psi(\mathbf r_1,\mathbf r_2,\ldots).
\]
Here $a$ labels the particles and $\boldsymbol\nabla_a$ differentiates with
respect to $\mathbf r_a$. The expression in parentheses
\SourcePage{63}{66}
is the operator of an infinitesimal displacement.

If a transformation does not change the Hamiltonian, transforming
$\hat H\psi$ gives the same result as first transforming $\psi$ and then
applying $\hat H$. If $\hat O$ produces the transformation, then
$\hat O(\hat H\psi)=\hat H(\hat O\psi)$, so
$\hat O\hat H-\hat H\hat O=0$: $\hat H$ commutes with $\hat O$.

Here $\hat O$ is the infinitesimal-displacement operator. The unit operator
commutes with every operator, and the constant $\delta\mathbf r$ may be taken
outside $\hat H$; hence
\begin{equation}
 \left(\sum_a\boldsymbol\nabla_a\right)\hat H
 -\hat H\left(\sum_a\boldsymbol\nabla_a\right)=0.
 \tag{15.1}
\end{equation}
An explicitly time-independent operator commuting with the Hamiltonian
corresponds to a conserved quantity. The quantity whose conservation follows
from the homogeneity of space is momentum (cf. Vol.~I, \S~7). Thus (15.1)
expresses momentum conservation; $\sum_a\boldsymbol\nabla_a$ corresponds, up
to a constant factor, to the total momentum, and each term to the momentum of
one particle.

The factor relating the momentum operator to $\boldsymbol\nabla$ follows from
the classical limit and equals $-i\hbar$:
\begin{equation}
 \hat{\mathbf p}=-i\hbar\boldsymbol\nabla,
 \tag{15.2}
\end{equation}
or
\[
 \hat p_x=-i\hbar\frac{\partial}{\partial x},\qquad
 \hat p_y=-i\hbar\frac{\partial}{\partial y},\qquad
 \hat p_z=-i\hbar\frac{\partial}{\partial z}.
\]
Indeed, using (6.1),
\[
 \hat{\mathbf p}\Psi=-i\hbar\boldsymbol\nabla\Psi
 =\Psi\boldsymbol\nabla S.
\]
\SourcePage{64}{67}
Thus classically the operator acts by multiplication by
$\boldsymbol\nabla S$, which is the classical momentum (see Vol.~I, \S~43).

The operator (15.2) is Hermitian, as required. For functions $\psi(x)$ and
$\varphi(x)$ vanishing at infinity, integration by parts gives
\[
 \int\varphi\hat p_x\psi\,dx
 =-i\hbar\int\varphi\frac{\partial\psi}{\partial x}\,dx
 =i\hbar\int\psi\frac{\partial\varphi}{\partial x}\,dx
 =\int\psi\hat p_x^*\varphi\,dx,
\]
which is the Hermiticity condition.

Since differentiations with respect to different variables commute,
\begin{equation}
 \hat p_x\hat p_y-\hat p_y\hat p_x=0,\quad
 \hat p_x\hat p_z-\hat p_z\hat p_x=0,\quad
 \hat p_y\hat p_z-\hat p_z\hat p_y=0.
 \tag{15.3}
\end{equation}
Thus all three momentum components may have definite values simultaneously.

Their eigenfunctions and eigenvalues satisfy
\begin{equation}
 -i\hbar\boldsymbol\nabla\psi=\mathbf p\psi,
 \tag{15.4}
\end{equation}
with solutions
\begin{equation}
 \psi=\mathrm{const}\cdot e^{i\mathbf p\cdot\mathbf r/\hbar}.
 \tag{15.5}
\end{equation}
The three components completely determine the wave function and form one
possible complete set. Their eigenvalues form a continuous spectrum from
$-\infty$ to $\infty$.

By (5.4), $\int\psi_{\mathbf p'}^*\psi_{\mathbf p},dV$ over all space should
equal $\delta(\mathbf p'-\mathbf p)$\footnote{For a vector argument,
$\delta(\mathbf a)=\delta(a_x)\delta(a_y)\delta(a_z)$.}. For later applications
it is more natural to normalize to the delta function of the momentum
difference divided by $2\pi\hbar$:
\[
 \int\psi_{\mathbf p'}^*\psi_{\mathbf p},dV
 =\delta\left(\frac{\mathbf p'-\mathbf p}{2\pi\hbar}\right),
\]
\SourcePage{65}{68}
or equivalently
\begin{equation}
 \int\psi_{\mathbf p'}^*\psi_{\mathbf p},dV
 =(2\pi\hbar)^3\delta(\mathbf p'-\mathbf p).
 \tag{15.6}
\end{equation}
Each one-dimensional factor obeys, for example,
$\delta[(p_x'-p_x)/(2\pi\hbar)]=2\pi\hbar\delta(p_x'-p_x)$.
Integration uses\footnote{The formula means that its left-hand side has the
defining property (5.8) of the delta function. Substitution in (5.8) gives the
Fourier formula
\[
 f(a)=\int_{-\infty}^{\infty}\int_{-\infty}^{\infty}
 f(x)e^{i\xi(x-a)}\,dx\,\frac{d\xi}{2\pi}.
\]}
\begin{equation}
 \frac{1}{2\pi}\int_{-\infty}^{\infty}e^{i\alpha\xi}\,d\xi=\delta(\alpha).
 \tag{15.7}
\end{equation}
It follows that $\mathrm{const}=1$ in (15.5)\footnote{With this normalization
$|\psi|^2=1$: the function is normalized to ``one particle per unit volume.''
This agreement is not accidental; cf. the note on p.~221.}:
\begin{equation}
 \psi_{\mathbf p}=e^{i\mathbf p\cdot\mathbf r/\hbar}.
 \tag{15.8}
\end{equation}

Expansion of an arbitrary $\psi(\mathbf r)$ in momentum eigenfunctions is its
Fourier integral:
\begin{equation}
 \psi(\mathbf r)=\int a(\mathbf p)\psi_{\mathbf p}(\mathbf r)
 \frac{d^3p}{(2\pi\hbar)^3}
 =\int a(\mathbf p)e^{i\mathbf p\cdot\mathbf r/\hbar}
 \frac{d^3p}{(2\pi\hbar)^3}.
 \tag{15.9}
\end{equation}
Here $d^3p=dp_xdp_ydp_z$, and
\begin{equation}
 a(\mathbf p)=\int\psi(\mathbf r)\psi_{\mathbf p}^*(\mathbf r)\,dV
 =\int\psi(\mathbf r)e^{-i\mathbf p\cdot\mathbf r/\hbar}\,dV.
 \tag{15.10}
\end{equation}
The function $a(\mathbf p)$ is the momentum-representation wave function;
$|a(\mathbf p)|^2d^3p/(2\pi\hbar)^3$ is the probability that the momentum lies
in $d^3p$.
\SourcePage{66}{69}
Just as $\hat{\mathbf p}$ represents momentum in coordinate space, introduce
$\hat{\mathbf r}$ for the coordinates in momentum space, defined by
\begin{equation}
 \bar{\mathbf r}=\int a^*(\mathbf p)\hat{\mathbf r}a(\mathbf p)
 \frac{d^3p}{(2\pi\hbar)^3}.
 \tag{15.11}
\end{equation}
But also $\bar{\mathbf r}=\int\psi^*\mathbf r\psi,dV$. Substituting (15.9)
and integrating by parts,
\[
 \mathbf r\psi(\mathbf r)
 =\int i\hbar e^{i\mathbf p\cdot\mathbf r/\hbar}
 \frac{\partial a(\mathbf p)}{\partial\mathbf p}
 \frac{d^3p}{(2\pi\hbar)^3},
\]
and hence
\[
 \bar{\mathbf r}=\int i\hbar a^*(\mathbf p)
 \frac{\partial a(\mathbf p)}{\partial\mathbf p}
 \frac{d^3p}{(2\pi\hbar)^3}.
\]
Comparison with (15.11) gives
\begin{equation}
 \hat{\mathbf r}=i\hbar\frac{\partial}{\partial\mathbf p}.
 \tag{15.12}
\end{equation}
In this representation momentum is simply multiplication by $\mathbf p$.

Finally, express the operator of translation through an arbitrary finite
distance $\mathbf a$ in terms of $\hat{\mathbf p}$. By definition,
$\hat T_{\mathbf a}\psi(\mathbf r)=\psi(\mathbf r+\mathbf a)$. Expanding
$\psi(\mathbf r+\mathbf a)$ in a Taylor series gives
\[
 \psi(\mathbf r+\mathbf a)=\psi(\mathbf r)
 +\mathbf a\cdot\frac{\partial\psi(\mathbf r)}{\partial\mathbf r}+\ldots
\]
or, on introducing $\hat{\mathbf p}=-i\hbar\boldsymbol\nabla$,
\[
 \psi(\mathbf r+\mathbf a)=\left[1+\frac{i}{\hbar}\mathbf a\cdot\hat{\mathbf p}
 +\frac12\left(\frac{i}{\hbar}\mathbf a\cdot\hat{\mathbf p}\right)^2
 +\ldots\right]\psi(\mathbf r).
\]
Thus
\begin{equation}
 \hat T_{\mathbf a}=\exp\left(\frac{i}{\hbar}
 \mathbf a\cdot\hat{\mathbf p}\right).
 \tag{15.13}
\end{equation}
This is the required \textit{finite-displacement operator}.
